\begin{examplebox}
	\textbf{\textcolor{CExample}{例 1（基础）}}
	\textbf{\textcolor{CExample}{题目：}} 已知 $f(x)=x^3$，$g(x)=x$，判断 $h(x)=f(x)+g(x)$ 的奇偶性。\textbf{\textcolor{CExample}{解题步骤：}}
	\begin{description}[style=nextline, leftmargin=2.8em, labelsep=0.5em, itemsep=0.45em]
		\item[\textbf{第一步（判断奇偶）：}]
			计算得 $f(-x)=(-x)^3=-x^3=-f(x)$，故 $f$ 为奇函数；$g(-x)=-x=-g(x)$，故 $g$ 为奇函数。
			\[
				f(-x)=-f(x),\quad g(-x)=-g(x).
			\]
			\par\textbf{理由：} 根据奇偶函数定义直接验证。
		\item[\textbf{第二步（应用规则）：}]
			由于 $f$ 与 $g$ 均为奇函数，根据奇$\pm$奇=奇，$h(x)=f(x)+g(x)$ 为奇函数。
			\[
				f(-x)=-f(x),\,g(-x)=-g(x)\;\Rightarrow\;(f+g)(-x)=-(f+g)(x).
			\]
			\par\textbf{理由：} 使用奇偶运算加减规则。
		\item[\textbf{第三步（验证结论）：}]
			直接计算 $h(-x)=(-x)^3+(-x)=-x^3-x=-(x^3+x)=-h(x)$，确认 $h$ 为奇函数。
			\[
				h(-x)=-h(x).
			\]
			\par\textbf{理由：} 代入计算验证。
	\end{description}
	\textbf{\textcolor{CExample}{关键结论：}}
	\textbf{$h(x)=x^3+x$ 为奇函数。}

	\medskip
	\textbf{\textcolor{CExample}{例 2（稍变形）}}
	\textbf{\textcolor{CExample}{题目：}} 已知 $f(x)=x^2$，$g(x)=\cos x$，判断 $\varphi(x)=f(x)g(x)$ 的奇偶性。\textbf{\textcolor{CExample}{解题步骤：}}
	\begin{description}[style=nextline, leftmargin=2.8em, labelsep=0.5em, itemsep=0.45em]
		\item[\textbf{第一步（判断奇偶）：}]
			$f(-x)=(-x)^2=x^2=f(x)$，所以 $f$ 为偶函数；$g(-x)=\cos(-x)=\cos x=g(x)$，所以 $g$ 为偶函数。
			\[
				f(-x)=f(x),\quad g(-x)=g(x).
			\]
			\par\textbf{理由：} 根据偶函数定义验证。
		\item[\textbf{第二步（应用规则）：}]
			由于 $f$ 与 $g$ 均为偶函数，根据偶$\times$偶=偶，$\varphi(x)=f(x)g(x)$ 为偶函数。
			\[
				\begin{aligned}
					f(-x)&=f(x),\; g(-x)=g(x)\;\Rightarrow\;(fg)(-x)\\&=f(x)g(x)\\&=(fg)(x).
			\end{aligned}
			\]
			\par\textbf{理由：} 使用奇偶运算乘法规则。
		\item[\textbf{第三步（验证结论）：}]
			直接计算 $\varphi(-x)=f(-x)g(-x)=f(x)g(x)=\varphi(x)$，确认 $\varphi$ 为偶函数。
			\[
				\varphi(-x)=\varphi(x).
			\]
			\par\textbf{理由：} 代入计算验证。
	\end{description}
	\textbf{\textcolor{CExample}{关键结论：}}
	\textbf{$x^2\cos x$ 为偶函数。}

	\medskip
	\textbf{\textcolor{CExample}{例 3（综合应用）}}
	\textbf{\textcolor{CExample}{题目：}} 设 $f(x)=x^3$（奇函数），$g(x)=x^2$（偶函数），$h(x)=f(x)g(x)$，$k(x)=h(x)+f(x)$，判断 $k(x)$ 的奇偶性。\textbf{\textcolor{CExample}{解题步骤：}}
	\begin{description}[style=nextline, leftmargin=2.8em, labelsep=0.5em, itemsep=0.45em]
		\item[\textbf{第一步（乘积奇偶）：}]
			$f$奇$g$偶$\Rightarrow$ $h$ 为奇函数。
			\[
				\begin{aligned}
					f(-x)&=-f(x),\; g(-x)=g(x)\;\Rightarrow\;(fg)(-x)\\&=-f(x)g(x)\\&=-(fg)(x).
			\end{aligned}
			\]
			\par\textbf{理由：} 应用乘法规则。
		\item[\textbf{第二步（加法奇偶）：}]
			$h$ 为奇函数，$f$ 为奇函数，奇+奇=奇，所以 $k$ 为奇函数。
			\[
				h(-x)=-h(x),\,f(-x)=-f(x)\;\Rightarrow\;(h+f)(-x)=-(h+f)(x).
			\]
			\par\textbf{理由：} 应用加法规则。
		\item[\textbf{第三步（直接验证）：}]
			计算 $k(-x)=h(-x)+f(-x)=(-h(x))+(-f(x))=-(h(x)+f(x))=-k(x)$，确认 $k$ 为奇函数。
			\[
				k(-x)=-k(x).
			\]
			\par\textbf{理由：} 代入计算验证。
	\end{description}
	\textbf{\textcolor{CExample}{关键结论：}}
	\textbf{$k(x)=x^5+x^3$ 为奇函数。}
\end{examplebox}
